\chapter{Botulinum Toxin Injection}

\section*{Overview}
Botulinum toxin injection is the delivery of botulinum neurotoxin into specific muscles or glands to treat conditions of excessive muscle contraction, glandular hypersecretion, pain, or overactive smooth muscle. The toxin produces a temporary, dose-dependent weakening of the injected target.

\section*{Alternative Names}
Botox injection; botulinum toxin type A; botulinum toxin type B; BTX-A; BTX-B

\section*{Principle}
Botulinum toxin blocks the release of acetylcholine at the neuromuscular junction by cleaving SNARE proteins, producing localized muscle paralysis. The effect develops over several days and wears off as new nerve terminals form over three to six months.

\section*{History}
Botulinum toxin type A was approved for strabismus and blepharospasm in the 1980s, and its therapeutic and cosmetic applications have since expanded dramatically.

\section*{Why This Test or Procedure Is Ordered}
Ordered for blepharospasm, cervical dystonia, spasticity from stroke or cerebral palsy, chronic migraine, hyperhidrosis, overactive bladder, strabismus, and other disorders of excessive muscle tone, movement, or secretion.

\section*{Is This a Primary or Secondary Test or Procedure}
A primary treatment for several focal dystonias and spasticity; it is adjunctive in conditions such as migraine and spasticity where oral therapies are also used.

\section*{How This Test or Procedure Is Performed}
The target muscles are identified by palpation, electromyography, or ultrasound guidance. A fine needle delivers small volumes of reconstituted toxin into the selected sites, with dosing adjusted to the muscle size and the intended effect.

\section*{What Preparation Is Needed}
Informed consent, review of concurrent aminoglycoside or neuromuscular-blocking medications, and identification of the injection sites are needed. No fasting is required for most indications.

\section*{Contraindications and Precautions}
Known hypersensitivity to botulinum toxin or albumin, infection at the injection site, and concurrent use of aminoglycoside antibiotics that may potentiate weakness are contraindications. Precautions include pregnancy and neuromuscular disorders such as myasthenia gravis, and care to avoid unintended spread to adjacent muscles.

\section*{Risks}
Risks include local pain, bruising, weakness of nearby muscles, dysphagia or respiratory compromise when treating cervical or laryngeal muscles, and rarely distant spread of toxin.

\section*{What Happens After the Test or Procedure}
The full effect develops over several days, and functional improvement is often assessed at two to four weeks. Relief lasts approximately three to six months, after which repeat injection can be scheduled.

\section*{Understanding Your Results}
Improvement in spasm, pain, or hyperhidrosis indicates a response, while insufficient benefit may prompt dose adjustment, change of target site, or switching between type A and type B preparations. Loss of effect over time may reflect antibody-mediated resistance.

\section*{Normal Results}
Marked reduction of the targeted involuntary contraction, pain, or sweating with preserved function of unaffected muscles indicates a successful treatment outcome.

\section*{Follow-up Tests and Procedures}
Outcomes are reviewed at follow-up visits and the dose is titrated for subsequent cycles. Physical or occupational therapy, oral spasmolytics, or surgical options may be added if response is inadequate.

\section*{References}
\begin{enumerate}
  \item MedlinePlus. Botulinum toxin injection. U.S. National Library of Medicine; 2025.
  \item Current Medical Diagnosis and Treatment. McGraw-Hill; 2025.
\end{enumerate}

\clearpage
